\documentclass[12pt,a4paper]{article}

% ----------- Encoding / language -----------
\usepackage[utf8]{inputenc}
\usepackage[T1]{fontenc}
\usepackage[english]{babel}

% ----------- Math / layout -----------
\usepackage{amsmath,amssymb}
\usepackage{geometry}
\usepackage{setspace}
\usepackage{lmodern}

% ----------- Figures / floats -----------
\usepackage{graphicx}
\usepackage{float}
\usepackage{caption}
\usepackage{fancyhdr}

% ----------- Code listings -----------
\usepackage{listings}
\usepackage{xcolor}
\usepackage{booktabs}
\usepackage{array}

\geometry{top=2cm,bottom=2cm,left=2.5cm,right=2.5cm}
\setstretch{1.3}
\pagestyle{fancy}
\fancyhf{}
\rfoot{\thepage}
\captionsetup[figure]{position=bottom,justification=centering}

\definecolor{codegray}{rgb}{0.95,0.95,0.95}
\definecolor{commentcolor}{rgb}{0.4,0.4,0.4}
\definecolor{keywordcolor}{rgb}{0.1,0.1,0.6}
\definecolor{stringcolor}{rgb}{0.2,0.5,0.2}

% Romanian letters inside listings (pdfLaTeX reads UTF-8 bytes, so map them explicitly)
\lstset{
  literate=
    {Ă}{{\u{A}}}1 {Â}{{\^A}}1 {Î}{{\^I}}1 {Ș}{{\c{S}}}1 {Ț}{{\c{T}}}1
    {ă}{{\u{a}}}1 {â}{{\^a}}1 {î}{{\^\i}}1 {ș}{{\c{s}}}1 {ț}{{\c{t}}}1
    {Ş}{{\c{S}}}1 {ş}{{\c{s}}}1 {Ţ}{{\c{T}}}1 {ţ}{{\c{t}}}1,
  language=Python,
  basicstyle=\ttfamily\small,
  backgroundcolor=\color{codegray},
  commentstyle=\color{commentcolor}\itshape,
  keywordstyle=\color{keywordcolor}\bfseries,
  stringstyle=\color{stringcolor},
  numbers=left,
  numberstyle=\tiny\color{gray},
  numbersep=8pt,
  frame=single,
  rulecolor=\color{black!30},
  showstringspaces=false,
  tabsize=4,
  breaklines=true,
  captionpos=b
}

\newcommand{\code}[1]{\texttt{\small #1}}

\begin{document}

% ===================== TITLE PAGE =====================
\begin{titlepage}
\begin{center}
  \textbf{MINISTRY OF EDUCATION AND RESEARCH}\\
  \textbf{OF THE REPUBLIC OF MOLDOVA}\\[0.5em]
  \textbf{TECHNICAL UNIVERSITY OF MOLDOVA}\\
  \textbf{Faculty of Computers, Informatics and Microelectronics}\\
  \textbf{Department of Software Engineering and Automatics}\\[2em]

  \vspace{1.5cm}
  \textbf{\Large Tacu Igor \textemdash{} FAF-242}\\[3em]

  \textbf{\Large REPORT}\\[0.5em]
  Laboratory Work No.\ 1\\[0.5em]
  \textbf{Cryptography}\\[0.5em]
  \textbf{The Caesar Cipher and the Caesar Cipher with a Permutation}\\[3em]

  \vfill
  \begin{flushright}
    Lecturer:\\
    Maia Zaica\\
    DISA, FCIM, UTM
  \end{flushright}
  \vfill

  \large \textbf{Chi\c{s}in\u{a}u \textemdash{} 2026}
\end{center}
\end{titlepage}

\tableofcontents
\newpage

% ===================== 1. THEORY =====================
\section{Short Theory}

\subsection{The Caesar cipher}

The Caesar cipher replaces every letter of the plaintext by the letter found $k$
positions further in the alphabet, wrapping around at the end. With an alphabet of
$n$ letters whose positions are numbered from $0$, the secret key is
$k \in \{1,2,\dots,n-1\}$ and the same $k$ is used for both operations:
\begin{align*}
  c &= e_k(x) = (x + k) \bmod n, \\
  m &= d_k(y) = (y - k) \bmod n.
\end{align*}
When $y-k$ is negative, $n$ must be added to bring the result back into the range
$0\dots n-1$, e.g.\ $(1-3) \bmod 26 = 24$. For the Romanian alphabet of this work
$n = 31$, so $k \in \{1,\dots,30\}$.

\textbf{Key space:} only $n-1 = 30$ keys. The cipher is broken by an
\emph{exhaustive search}: the attacker decrypts with every key and keeps the one
that yields readable text.

\subsection{The Caesar cipher with a permutation}

To enlarge the key space, the alphabet is first permuted using a \emph{keyword}
$k_2$: the distinct letters of $k_2$ are written first (only the first occurrence of
each is kept), followed by the remaining letters of the alphabet in their natural
order. The Caesar shift $k_1$ is then applied to the \emph{positions in the permuted
alphabet}. For $k_2 = \textit{cryptography}$ and the English alphabet the order becomes
\texttt{CRYPTOGAHBDEFIJKLMNQSUVWXZ}, so with $k_1 = 3$ the text \texttt{CIFRULCEZAR}
is encrypted as \texttt{PLKTXQPJYDT} (instead of \texttt{FLIUXOFHCDU} with the plain
cipher). Encryption and decryption keep the same formulas, but $x$ and $y$ are now
positions in the permuted alphabet.

\textbf{Key space:} there are $n!$ orderings and $n-1$ shifts. For $n = 31$:
\[
  31! \times 30 \approx 8.22\cdot 10^{33} \times 30 \approx 2.47\cdot 10^{35}
  \quad (\approx 2^{117.6}),
\]
which makes exhaustive search infeasible. However, a keyword always yields a
\emph{fixed} substitution, so letter frequencies of the language remain visible in
the ciphertext and the cipher is still vulnerable to \emph{frequency analysis}. The
number of distinct encryptions cannot exceed $n!$, because different key pairs can
produce the same substitution.

% ===================== 2. IMPLEMENTATION =====================
\section{Implementation}

The program is written in Python 3 (no external libraries) and consists of
\code{caesar.py} (cipher + console interface) and \code{test\_caesar.py} (unit
tests). The alphabet is a parameter of every function, which allowed me to verify the
code against the English examples from the assignment.

\subsection{Letter encoding}

\begin{lstlisting}
ROMANIAN = ["A", "Ă", "Â", "B", "C", "D", "E", "F", "G", "H", "I", "Î", "J", "K", "L",
            "M", "N", "O", "P", "Q", "R", "S", "Ș", "T", "Ț", "U", "V", "W", "X", "Y", "Z"]
CEDILLA = {"Ş": "Ș", "ş": "ș", "Ţ": "Ț", "ţ": "ț"}
\end{lstlisting}
The code of a letter is simply its index in this list, which reproduces Table~2 of
the assignment ($\text{A}=0$, $\text{\u{A}}=1$, \dots, $\text{Z}=30$). Built-in
character codes (\code{ord}, \code{chr}, ASCII/Unicode) are \emph{never} used to
compute the shift. \code{CEDILLA} maps the cedilla letters \c{S}, \c{T} to the
comma-below letters so that text typed on different keyboards is accepted.

\subsection{Input validation}

\begin{lstlisting}
def normalize(text, alphabet=ROMANIAN):
    out = ""
    for ch in text.replace(" ", ""):
        ch = CEDILLA.get(ch, ch).upper()
        if ch not in alphabet:
            raise InvalidInput(f"Invalid character '{ch}'. Allowed: ...")
        out += ch
    return out

def validate_shift(raw, n=len(ROMANIAN)):
    try:
        k = int(raw)
    except ValueError:
        k = None
    if k is None or not 1 <= k < n:
        raise InvalidInput(f"Invalid key '{raw.strip()}'. The key must be an "
                           f"integer between 1 and {n - 1} inclusive.")
    return k
\end{lstlisting}
\code{normalize} removes spaces, unifies cedilla variants, converts to uppercase
(requirement~6) and rejects the first character that is not in the alphabet,
naming it in the message. \code{validate\_shift} accepts only integers in
$1\dots30$. \code{validate\_keyword} (not shown) applies \code{normalize} to the
keyword and additionally rejects keywords with spaces or shorter than $7$ letters.
All validators raise \code{InvalidInput}; the function \code{ask} catches it, prints
the message and asks for the value again:
\begin{lstlisting}
def ask(prompt, validate):
    while True:
        try:
            return validate(input(prompt))
        except InvalidInput as err:
            print(f"  [!] {err}")
\end{lstlisting}

\subsection{Encryption and decryption (Task 1.1)}

\begin{lstlisting}
def shift(text, k, alphabet):
    n = len(alphabet)
    pos = {ch: i for i, ch in enumerate(alphabet)}
    return "".join(alphabet[(pos[ch] + k) % n] for ch in text)

def encrypt(text, k, alphabet=ROMANIAN):   # c = (x + k) mod n
    return shift(text, k, alphabet)

def decrypt(text, k, alphabet=ROMANIAN):   # m = (y - k) mod n
    return shift(text, -k, alphabet)
\end{lstlisting}
\code{pos} maps each letter to its position and \code{shift} applies the formulas of
Section~1; decryption is the same operation with $-k$. Python's \code{\%} always returns a
value in $0\dots n-1$, so a negative $y-k$ is wrapped around correctly (e.g.\ $(1-3) \bmod 31 = 29$).

\subsection{Permuted alphabet and two keys (Task 1.2)}

\begin{lstlisting}
def permuted_alphabet(keyword, alphabet=ROMANIAN):
    # dict keeps insertion order: repeats dropped, first occurrence kept
    return list(dict.fromkeys(keyword + "".join(alphabet)))

def encrypt2(text, k, keyword, alphabet=ROMANIAN):
    return encrypt(text, k, permuted_alphabet(keyword, alphabet))
\end{lstlisting}
Putting the keyword in front of the whole alphabet and removing repeated letters, produces exactly the construction from the assignment (first
occurrence of each keyword letter, then the remaining letters in natural order).
The resulting order simply replaces the alphabet in \code{encrypt}/\code{decrypt},
so the second task reuses the code of the first. The program prints the permuted
alphabet before every operation so that the result can be checked by hand.

\subsection{Verification}

\code{test\_caesar.py} checks the code against all worked examples of the assignment
(using the English alphabet): $\texttt{CIFRULCEZAR}\to\texttt{FLIUXOFHCDU}$
($k=3$), $\texttt{BRUTEFORCEATTACK}\to\texttt{SILKVWFITVRKKRTB}$ ($k=17$) together
with the brute-force candidate for $k=1$, the permuted alphabet for
\textit{cryptography}, and $\texttt{CIFRULCEZAR}\to\texttt{PLKTXQPJYDT}$ with
$k_1=3$, $k_2=\textit{cryptography}$. For the Romanian alphabet it checks the size
(31) and codes, the wrap-around cases, round trips for all $30$ keys, and the
rejected inputs. All 9 tests pass (\code{python3 -m unittest}).

% ===================== 3. RESULTS =====================
\section{Results}

The listings below are the real output of the program (the text typed by the user
is shown after each prompt).

\subsection{Task 1.1 -- Caesar cipher}

\begin{lstlisting}[language={},numbers=none,basicstyle=\ttfamily\footnotesize,caption={Encryption and decryption with $k=3$.}]
Choice: 1
Operation - (e)ncrypt or (d)ecrypt: e
Key 1 (shift, 1-30): 3
Enter the plaintext: cifrul cezar
Normalized input : CIFRULCEZAR
Ciphertext       : FKITXOFHÂBT

Choice: 1
Operation - (e)ncrypt or (d)ecrypt: d
Key 1 (shift, 1-30): 3
Enter the ciphertext: FKITXOFHÂBT
Normalized input : FKITXOFHÂBT
Plaintext        : CIFRULCEZAR
\end{lstlisting}
The result differs from \texttt{FLIUXOFHCDU} of the English example because the
Romanian alphabet has $31$ letters including \u{A}, \^A, \^I, \c{S}, \c{T}: for
example $\text{C}=4 \to 7=\text{F}$ and $\text{I}=10\to 13=\text{K}$.

\begin{lstlisting}[language={},numbers=none,basicstyle=\ttfamily\footnotesize,caption={Task 1.1 -- invalid key and invalid character.}]
Key 1 (shift, 1-30): 0
  [!] Invalid key '0'. The key must be an integer between 1 and 30 inclusive.
Key 1 (shift, 1-30): 31
  [!] Invalid key '31'. The key must be an integer between 1 and 30 inclusive.
Key 1 (shift, 1-30): abc
  [!] Invalid key 'abc'. The key must be an integer between 1 and 30 inclusive.
Key 1 (shift, 1-30): 5
Enter the plaintext: hello, world
  [!] Invalid character ','. Allowed: letters of the Romanian alphabet (A-Z, Ă, Â, Î, Ș, Ț) and spaces.
Enter the plaintext: Ştiinţă şi artă
Normalized input : ȘTIINȚĂȘIARTĂ
Ciphertext       : WXMMSYEWMDUXE
\end{lstlisting}
The last line shows the cedilla letters being accepted as \c{S} and \c{T}.

\subsection{Task 1.2 -- Caesar cipher with two keys}

\begin{lstlisting}[language={},numbers=none,basicstyle=\ttfamily\footnotesize,caption={Encryption and decryption with $k_1=3$, $k_2=\textit{criptografie}$.}]
Choice: 2
Operation - (e)ncrypt or (d)ecrypt: e
Key 1 (shift, 1-30): 3
Key 2 (keyword, min 7 letters): criptografie
Enter the plaintext: Cifrul Cezar
Permuted alphabet: C R I P T O G A F E Ă Â B D H Î J K L M N Q S Ș Ț U V W X Y Z
Normalized input : CIFRULCEZAR
Ciphertext       : POÂTXQPBIĂT

Choice: 2
Operation - (e)ncrypt or (d)ecrypt: d
Key 1 (shift, 1-30): 3
Key 2 (keyword, min 7 letters): criptografie
Enter the ciphertext: POÂTXQPBIĂT
Permuted alphabet: C R I P T O G A F E Ă Â B D H Î J K L M N Q S Ș Ț U V W X Y Z
Normalized input : POÂTXQPBIĂT
Plaintext        : CIFRULCEZAR
\end{lstlisting}
Check by hand: C has position $0$ in the permuted alphabet, so it becomes the letter
at position $3$, which is P; I (position $2$) becomes the letter at position $5$, O.

\begin{lstlisting}[language={},numbers=none,basicstyle=\ttfamily\footnotesize,caption={Task 1.2 -- invalid keywords.}]
Key 2 (keyword, min 7 letters): cripto
  [!] Keyword too short (6 letters). It needs at least 7.
Key 2 (keyword, min 7 letters): cripto grafie
  [!] Invalid keyword. It must be a single word without spaces.
Key 2 (keyword, min 7 letters): cripto9rafie
  [!] Invalid character '9'. Allowed: letters of the Romanian alphabet (A-Z, Ă, Â, Î, Ș, Ț) and spaces.
Key 2 (keyword, min 7 letters): criptografie
\end{lstlisting}

% ===================== 4. CONCLUSIONS =====================
\section{Conclusions}

The plain Caesar cipher has only $n-1 = 30$ keys for the Romanian alphabet, so an
exhaustive search finds the plaintext immediately; it offers no real protection.
Using a keyword to permute the alphabet raises the number of keys to about
$31!\cdot 30 \approx 2.5\cdot10^{35}$ and makes brute force infeasible, but the
cipher remains a monoalphabetic substitution: every plaintext letter is always
replaced by the same ciphertext letter, so a frequency analysis of a sufficiently
long ciphertext still reveals the substitution. A large key space is therefore
necessary but not sufficient for security.

During the work I learned to implement modular arithmetic with explicit handling
of negative values, to build the permuted alphabet from a keyword, and to
validate user input with clear error messages. Two practical details mattered: the
Romanian alphabet must be encoded from the given table (not from Unicode code
points, which are not contiguous for \u{A}, \^A, \^I, \c{S}, \c{T}), and the
cedilla/comma-below variants of \c{S} and \c{T} must be unified. Testing the
alphabet-independent functions against the English examples of the assignment gave
confidence that the Romanian version is correct.

% ===================== 5. SOURCE CODE =====================
\section{Source Code}

The full source code is available in the GitHub repository:\\
\texttt{https://github.com/igortacu/criptography}

\end{document}
